\subsection{Paper 35: A Transparent Search-Based Framework for Skin and Breast Cancer Diagnosis Using Siamese Networks}
\label{sec:paper_35}

\begin{tcolorbox}[colback=gray!5!white,colframe=gray!75!black,title=\textbf{Metadata \& Citation Summary}]
\textbf{Authors:} Z. Tabatabaei; A. Colomer; J. O. Moll; V. Naranjo \\
\textbf{Journal / Venue:} \textit{IEEE Access} (2025) \\
\textbf{DOI:} \href{https://doi.org/10.1109/ACCESS.2025.3628157}{10.1109/ACCESS.2025.3628157} \quad | \quad \textbf{IEEE Article Number:} 11224374 \\
\textbf{Thematic Category:} \textsc{Siamese Metric Retrieval \& Explainable Diagnosis} \quad | \quad \textbf{Citation:} \cite{11224374}
\end{tcolorbox}

\subsubsection*{Research Context \& Problem Formulation}
Ambiguous melanocytic lesions (e.g. Spitzoid neoplasms) often cannot be definitively classified by binary labels alone, requiring dermatologists to review visually and histologically similar confirmed cases.

\paragraph{Literature Synthesis:} Histopathological diagnosis requires pathologists to interpret complex biopsy slides, where black-box AI classifiers provide diagnoses without interpretable reference precedents. To provide transparent decision support, the authors propose a Content-Based Histopathological Image Retrieval (CBHIR) framework evaluated on skin cancer histopathology patch collections and the benchmark BreakHis breast cancer dataset (9,109 microscopic images). Preprocessing incorporated stain normalization using the Macenko method to mitigate inter-laboratory staining variations, patch extraction, and data augmentation including rotations and elastic deformations. The proposed framework employs a deep Siamese Neural Network trained with contrastive and triplet loss functions to learn an invariant metric embedding space where histopathologically correlated lesions cluster tightly together. When a query patch is submitted, the system transparently retrieves the top-k most morphologically similar historical cases with verified diagnoses. The framework attained a Top-1 retrieval diagnosis accuracy of 88.50\%, a Top-5 accuracy of 94.20\%, a precision of 89.10\%, a recall/sensitivity of 87.80\%, and an F1-score of 88.40\%. Standalone AUC was not explicitly reported in the text. Key strengths include case-based interpretability that mirrors pathologist clinical workflows. Explicitly stated limitations include high memory and indexing latency when scaling to whole slide image databases, reliance on biopsy tissue rather than non-invasive imaging, and retrieval quality bound to reference archive diversity.

\subsubsection*{Formalized Research Questions (RQs)}
\begin{enumerate}[label=\textbf{\arabic*.}, leftmargin=*]
  \item \textbf{RQ1 [Content-Based Twin Retrieval]:} \textit{Can a Siamese Neural Network framework for Content-Based Image Retrieval (CBIR) provide transparent diagnostic decision support by retrieving visually and pathologically verified 'twin' cases?}
  \begin{itemize}[leftmargin=1.5em]
    \item \textbf{Investigation Methodology:} Trained a Siamese architecture with contrastive metric loss to map dermoscopic and histopathological images into an embedding space where pathologically identical tumors cluster together.
    \item \textbf{Empirical Answer \& Evidence:} Yes. The framework achieved a Top-5 retrieval precision of 94.2\% and classification accuracy of 93.8\%, retrieving verifiable twin cases with known histological outcomes.
  \end{itemize}
  \item \textbf{RQ2 [Contrastive Metric Learning Space]:} \textit{How effectively does contrastive metric learning separate borderline Spitzoid tumors from malignant melanomas compared to standard softmax classifiers?}
  \begin{itemize}[leftmargin=1.5em]
    \item \textbf{Investigation Methodology:} Analyzed t-SNE projections of learned Siamese metric embeddings on ambiguous Spitzoid tumors, conventional nevi, and melanomas.
    \item \textbf{Empirical Answer \& Evidence:} Contrastive metric learning formed clear semantic sub-clusters for borderline cases, avoiding the overconfident misclassifications typical of standard softmax heads.
  \end{itemize}
  \item \textbf{RQ3 [Cross-Organ Generalization]:} \textit{Does the transparent search-based Siamese retrieval framework generalize across diverse oncology domains (skin cancer and breast cancer)?}
  \begin{itemize}[leftmargin=1.5em]
    \item \textbf{Investigation Methodology:} Evaluated identical Siamese metric learning pipelines on dermoscopic melanoma datasets and histopathological breast cancer archives.
    \item \textbf{Empirical Answer \& Evidence:} Demonstrated robust cross-domain generalizability, achieving high retrieval fidelity and diagnostic concordance across both dermatological and breast cancer benchmarks.
  \end{itemize}
  \item \textbf{RQ4 [Large-Scale Retrieval Scalability]:} \textit{What is the search latency overhead of querying massive multi-thousand case clinical repositories during point-of-care consultations?}
  \begin{itemize}[leftmargin=1.5em]
    \item \textbf{Investigation Methodology:} Benchmarked exact nearest-neighbor search latency against approximate indexing algorithms (FAISS, HNSW) across scaling database sizes.
    \item \textbf{Empirical Answer \& Evidence:} Exact nearest-neighbor search scales linearly with repository size, requiring approximate nearest-neighbor indexing (e.g., HNSW) to maintain sub-second retrieval in massive archives.
  \end{itemize}
\end{enumerate}

\subsubsection*{Experimental Benchmark \& Quantitative Highlights}
\begin{itemize}[leftmargin=*]
  \item \textbf{Benchmark Repositories:} Clinical Spitzoid Neoplasm Cohort, ISIC Melanoma Benchmark, Breast Cancer Histopathology Archive
  \item \textbf{Key Empirical Results:} \textbf{Top-5 Retrieval Precision}: 94.2\% \quad | \quad \textbf{Diagnostic Accuracy}: 93.8\% \quad | \quad \textbf{Multi-Organ Compatibility}: Validated on Skin \& Breast Cancer
\end{itemize}

\subsubsection*{Open Research Questions \& Future Directions}
\begin{itemize}[leftmargin=*]
  \item \textit{Can generative text models automatically synthesize a comparative clinical summary contrasting the query image with retrieved twin cases?}
  \item \textit{How does the retrieval system handle domain shift when query images are captured using different dermatoscope optical models?}
\end{itemize}
\vspace{0.8em}\hrule\vspace{1.2em}
